\chapter{Reproduction Instructions\ifFullReport{} and Command Ledger\fi}\label{app:reproduction}
\section{Environment and setup}
All commands below are issued from the delivered project root unless a directory change is shown. The setup script selects a supported CPython interpreter, creates a local virtual environment, installs the pinned requirements and local package, obtains the official learned-matcher source at its recorded revision, and downloads the pretrained checkpoints. The data directory, model directory and environment are excluded from source control.
\begin{lstlisting}[numbers=none]
bash scripts/setup.sh
make data
make baseline
make evaluate
make train
make export
make demo
make profile
make test
make thesis
\end{lstlisting}
The trained-head and export steps use the frozen data preparation. The demonstration and artifact generator consume the main evaluation results. The complete-query profile is executed after other heavy inference tasks finish. The output PDF is \code{report/thesis.pdf}; comparison clips and animation are in \code{demo/}. Exact arguments used in this run are preserved in \code{results/commands.jsonl}\ifFullReport{} and in the ledger later in this appendix\fi.

\section{Direct CLI checks}
\begin{lstlisting}[numbers=none]
.venv/bin/python scripts/run_logged.py .venv/bin/stitch index data/places.json
.venv/bin/python scripts/run_logged.py .venv/bin/stitch query data/places.json test.jpg results/raw/cli_query.png
.venv/bin/python scripts/run_logged.py .venv/bin/stitch video data/places.json video.mp4 results/raw/cli_video.mp4 --limit 8
.venv/bin/python scripts/run_logged.py .venv/bin/python scripts/download_data.py --verify-only
\end{lstlisting}
The query image is already represented in this demonstration database, so this CLI check is a functional content-retrieval check rather than a held-out accuracy result. The video command exercises decoding, EOF/limit handling, recognition and writing. Bounded file inputs replace the unavailable private DroidCam stream; no live phone capture is claimed.

\section{Data integrity and restoration}
\code{results/data_manifest.json} pins downloaded filenames, lengths and SHA256 hashes. \code{results/preparation.json} pins generated synthetic inputs and the frozen benchmark manifest. Existing valid collections are verified and skipped by the downloader. Missing collections can be restored from the recorded public sources and checked against the supplied manifest. A corrupt existing file fails explicitly, preventing unnoticed input replacement.

The original HPatches host and author mirror are documented in \cref{ch:datasets}. Tokyo subset members are verified against ZIP CRC and decompressed size before their local hashes are recorded. Full-archive transfer failures remain in the command journal. Reproduction depends on the continued availability and terms of the original sources; their presence in this artifact is not a promise that a remote URL will remain unchanged.

\ifFullReport
\section{Manual datasets for expanded evaluation}
\subsection{Aachen Day-Night v1.1}
The official landing page is \url{https://www.visuallocalization.net/datasets/}. The inspected public archive root is \url{https://data.ciirc.cvut.cz/public/projects/2020VisualLocalization/Aachen-Day-Night/}. Follow both official README files, obtaining base images, the v1.1 extension, query intrinsics and the appropriate 3D reconstruction. The extension is not a substitute for every required base file. A proposed local layout is:
\begin{lstlisting}[numbers=none]
data/manual/aachen/
  images/
    database/
    query/
  queries/
  3D-models/
  aachen_v1_1/
\end{lstlisting}
Actual names should follow the archive README; an adapter should map them into explicit image, camera and 3D-point records. Query-pose evaluation may require the benchmark's evaluation procedure rather than public ground truth. The official image/model terms specify non-commercial academic restrictions. No Aachen pose score is supplied by the current artifact.

\subsection{RobotCar Seasons}
Use \url{https://data.ciirc.cvut.cz/public/projects/2020VisualLocalization/RobotCar-Seasons/}. Obtain the condition image archives, camera intrinsics/extrinsics, selected model and the correct version's split lists. A proposed adapter root is:
\begin{lstlisting}[numbers=none]
data/manual/robotcar/
  images/<condition>/<camera>/
  intrinsics/
  extrinsics/
  3D-models/
  robotcar_v2_train.txt
  robotcar_v2_test.txt
\end{lstlisting}
The saved official documentation describes the licensing and version distinctions. These images were not downloaded for the executed experiment. Their night and seasonal conditions could extend the present evidence, but must not be reported until acquisition, calibration and evaluation are actually run.

\subsection{Extended CMU Seasons}
Use \url{https://data.ciirc.cvut.cz/public/projects/2020VisualLocalization/Extended-CMU-Seasons/}. Each slice archive contains the corresponding images and map-related data; choose slices and conditions before inference and retain that selection. An expected adapter structure is:
\begin{lstlisting}[numbers=none]
data/manual/cmu/
  intrinsics.txt
  slice2/
    database/
    query/
    sparse/
    camera_poses/
\end{lstlisting}
Archive-internal directory names should be taken from the official README. The inspected source specifies attribution, non-commercial and share-alike restrictions. No CMU image or seasonal metric is included in this thesis.

\subsection{MSLS and the full Tokyo street-view task}
MSLS instructions are at \url{https://github.com/mapillary/mapillary_sls} and its linked download site \url{https://www.mapillary.com/dataset/places}. Preserve city, query/database and sequence metadata. The official structure uses \code{train_val/<city>/query/images/}, a corresponding database directory and associated CSV files. Test data and current ground-truth availability should follow the current official instructions rather than an assumed historical evaluation server.

For the official Tokyo street-view experiment, follow \url{https://github.com/gmberton/VPR-datasets-downloader/blob/main/format_tokyo247.py} and the original dataset's access instructions. That formatting source explains that the database must be obtained separately; the publicly downloaded query archive alone does not reproduce the official benchmark. A suitable root retains street-view panorama crops, the dataset metadata and original queries in distinct directories.

Pittsburgh/SF-XL, Oxford/Paris and Google Landmarks are alternative future collections, not failed downloads in this run. They should be documented with their own licences and split protocols if used later. None is needed to reproduce the results already presented.
\fi

\section{Build procedure}
The installed TeX toolchain was used directly. The document uses the report class, numeric natbib citations, generated vector figures and a nomenclature index. From the project root, \code{make thesis} regenerates empirical artifacts and invokes latexmk. The thesis Makefile copies the completed PDF to the main deliverable path and runs the verification script. A clean TeX rebuild is available through the recorded clean/build commands.

Authoring and environment-inspection commands performed before the command runner existed are summarised in \code{results/bootstrap_commands.json}. Subsequent substantive executions are recorded in \code{results/commands.jsonl} with exact command text and outcomes\ifFullReport{}, reproduced below\fi. File edits are represented by the delivered source, not presented as experiment commands. \ifFullReport The currently running build can appear as in progress when its own ledger is generated; the machine-readable journal records completion after the build returns.\fi

\ifFullReport
\section{Exact recorded execution ledger}
\noindent\textbf{Initial setup/inspection 1}.
\begin{lstlisting}[numbers=none]
pwd && rg --files -g 'AGENTS.md' -g '!node_modules' -g '!venv' -g '!\.git' . /Users/bhattacharjee/.codex/attachments/96f3cc0c-84ef-4a82-a64b-b9fce597c78b 2>/dev/null
\end{lstlisting}

\noindent\textbf{Initial setup/inspection 2}.
\begin{lstlisting}[numbers=none]
cat '/Users/bhattacharjee/.codex/attachments/96f3cc0c-84ef-4a82-a64b-b9fce597c78b/Pasted text.txt'
\end{lstlisting}

\noindent\textbf{Initial setup/inspection 3}.
\begin{lstlisting}[numbers=none]
ls -la
\end{lstlisting}

\noindent\textbf{Initial setup/inspection 4}.
\begin{lstlisting}[numbers=none]
for p in /AGENTS.md /Users/AGENTS.md /Users/bhattacharjee/AGENTS.md /Users/bhattacharjee/Desktop/AGENTS.md; do if [ -f "$p" ]; then cat "$p"; fi; done
find . -name AGENTS.md -print
\end{lstlisting}

\noindent\textbf{Initial setup/inspection 5}.
\begin{lstlisting}[numbers=none]
python3 --version
which git pdflatex latexmk biber brew ffmpeg tectonic
 df -h .
sysctl -n machdep.cpu.brand_string hw.memsize
\end{lstlisting}

\noindent\textbf{Initial setup/inspection 6}.
\begin{lstlisting}[numbers=none]
git clone https://github.com/BitsJayMehta173/Stitch.git Stitch-research
\end{lstlisting}

\noindent\textbf{Initial setup/inspection 7}.
\begin{lstlisting}[numbers=none]
ls /opt/homebrew/bin/python* /usr/local/bin/python* /usr/bin/python* 2>/dev/null
uname -m
system_profiler SPHardwareDataType | head -22
kpsewhich algorithm.sty
kpsewhich algpseudocode.sty
kpsewhich siunitx.sty
\end{lstlisting}

\noindent\textbf{Initial setup/inspection 8}.
\begin{lstlisting}[numbers=none]
find /opt/homebrew/bin /usr/local/bin /usr/bin -maxdepth 1 -name 'python*' -print 2>/dev/null
\end{lstlisting}

\noindent\textbf{Initial setup/inspection 9}.
\begin{lstlisting}[numbers=none]
/usr/bin/python3 --version
/opt/homebrew/bin/python3.14 --version
ls /opt/homebrew/opt | rg 'python|uv'
find /Users/bhattacharjee/Desktop/internship/Stitch-research -maxdepth 2 -type f -not -path '*/.git/*' -print
\end{lstlisting}

\noindent\textbf{Initial setup/inspection 10}.
\begin{lstlisting}[numbers=none]
ls -la Stitch-research
curl -I --connect-timeout 10 --max-time 20 https://raw.githubusercontent.com/BitsJayMehta173/Stitch/main/tep3.py
\end{lstlisting}

\noindent\textbf{Initial setup/inspection 11}.
\begin{lstlisting}[numbers=none]
git rev-parse HEAD
rg --files
cat tep.py tep2.py tep3.py README.md
\end{lstlisting}

\noindent\textbf{Initial setup/inspection 12}.
\begin{lstlisting}[numbers=none]
/usr/bin/python3 -m venv .venv
mkdir -p src/stitch scripts configs tests data results/raw results/logs figures thesis/chapters thesis/figures thesis/tables demo vendor
.venv/bin/python -m pip --version
\end{lstlisting}

\noindent\textbf{20261105T051728.588219Z}: exit 0, 2.263 seconds.
\begin{lstlisting}[numbers=none]
.venv/bin/python -m pip install --upgrade pip
\end{lstlisting}

\noindent\textbf{20261105T051730.887161Z}: exit 0, 150.577 seconds.
\begin{lstlisting}[numbers=none]
.venv/bin/python -m pip install -r requirements.txt
\end{lstlisting}

\noindent\textbf{20261105T051803.763940Z}: exit 0, 14.033 seconds.
\begin{lstlisting}[numbers=none]
git clone --depth 1 https://github.com/cvg/LightGlue.git vendor/LightGlue
\end{lstlisting}

\noindent\textbf{20261105T051817.840310Z}: exit 0, 1.409 seconds.
\begin{lstlisting}[numbers=none]
git clone --depth 1 https://github.com/gmberton/CosPlace.git vendor/CosPlace
\end{lstlisting}

\noindent\textbf{20261105T051830.718242Z}: exit 6, 1.696 seconds.
\begin{lstlisting}[numbers=none]
curl -IL --connect-timeout 10 --max-time 25 https://icvl.ee.ic.ac.uk/vbalnt/hpatches/hpatches-sequences-release.tar.gz
\end{lstlisting}

\noindent\textbf{20261105T051832.447559Z}: exit 0, 3.881 seconds.
\begin{lstlisting}[numbers=none]
curl -IL --connect-timeout 10 --max-time 25 'http://www.ok.sc.e.titech.ac.jp/~torii/project/247/'
\end{lstlisting}

\noindent\textbf{20261105T051859.042356Z}: exit 28, 180.034 seconds.
\begin{lstlisting}[numbers=none]
curl -fL --connect-timeout 10 --max-time 180 -o data/tokyo_queries.zip https://data.ciirc.cvut.cz/public/projects/2015netVLAD/Tokyo247/queries/247query_subset_v2.zip
\end{lstlisting}

\noindent\textbf{20261105T051946.478990Z}: exit 0, 479.281 seconds.
\begin{lstlisting}[numbers=none]
curl -fL --connect-timeout 10 --max-time 1200 -o data/hpatches-sequences-release.zip https://huggingface.co/datasets/vbalnt/hpatches/resolve/main/hpatches-sequences-release.zip
\end{lstlisting}

\noindent\textbf{20261105T052006.808242Z}: exit 0, 1.099 seconds.
\begin{lstlisting}[numbers=none]
curl -sSIL --connect-timeout 10 --max-time 20 https://huggingface.co/datasets/vbalnt/hpatches/resolve/main/hpatches-sequences-release.zip
\end{lstlisting}

\noindent\textbf{20261105T052007.946358Z}: exit 0, 0.845 seconds.
\begin{lstlisting}[numbers=none]
curl -sS -I -H 'Range: bytes=0-99' --connect-timeout 10 --max-time 20 https://data.ciirc.cvut.cz/public/projects/2015netVLAD/Tokyo247/queries/247query_subset_v2.zip
\end{lstlisting}

\noindent\textbf{20261105T052041.129322Z}: exit 1, 40.940 seconds.
\begin{lstlisting}[numbers=none]
.venv/bin/python -c 'import sys,zipfile;sys.path.insert(0,'"'"'scripts'"'"');from remote_zip import RemoteZipIO;r=RemoteZipIO('"'"'https://data.ciirc.cvut.cz/public/projects/2015netVLAD/Tokyo247/queries/247query_subset_v2.zip'"'"');z=zipfile.ZipFile(r);print(len(z.namelist()));print('"'"'\n'"'"'.join(z.namelist()[:35]));print([(n,z.read(n)[:150]) for n in z.namelist() if n.endswith('"'"'.csv'"'"')][:1])'
\end{lstlisting}

\noindent\textbf{20261105T052159.098463Z}: exit 0, 2.430 seconds.
\begin{lstlisting}[numbers=none]
curl -fsSL --connect-timeout 10 --max-time 30 -o data/hpatches_hf_tree.json 'https://huggingface.co/api/datasets/JimmyFu/hpatches_sequences/tree/main?recursive=true'
\end{lstlisting}

\noindent\textbf{20261105T052302.124472Z}: exit 0, 15.673 seconds.
\begin{lstlisting}[numbers=none]
.venv/bin/python -c 'import sys,zipfile,json;sys.path.insert(0,'"'"'scripts'"'"');from remote_zip import RemoteZipIO;r=RemoteZipIO('"'"'https://data.ciirc.cvut.cz/public/projects/2015netVLAD/Tokyo247/queries/247query_subset_v2.zip'"'"');z=zipfile.ZipFile(r);print(z.read('"'"'247query_subset_v2/Readme.txt'"'"').decode());n=next(n for n in z.namelist() if n.endswith('"'"'.csv'"'"'));print(n,z.read(n).decode());open('"'"'data/tokyo_zip_directory.json'"'"','"'"'w'"'"').write(json.dumps([dict(name=i.filename,offset=i.header_offset,size=i.file_size,compressed=i.compress_size,crc=i.CRC,method=i.compress_type) for i in z.infolist()],indent=2))'
\end{lstlisting}

\noindent\textbf{20261105T052305.737685Z}: exit 0, 129.830 seconds.
\begin{lstlisting}[numbers=none]
curl -fsSL --connect-timeout 10 --max-time 180 -o models_cosplace_download.pth https://github.com/gmberton/CosPlace/releases/download/v1.0/ResNet18_512_cosplace.pth
\end{lstlisting}

\noindent\textbf{20261105T052317.888293Z}: exit 0, 2.865 seconds.
\begin{lstlisting}[numbers=none]
.venv/bin/python -m pip install -e . --no-deps
\end{lstlisting}

\noindent\textbf{20261105T052320.783618Z}: exit 0, 2.437 seconds.
\begin{lstlisting}[numbers=none]
.venv/bin/python -m pip install -e vendor/LightGlue --no-deps
\end{lstlisting}

\noindent\textbf{20261105T052454.115168Z}: exit 0, 386.097 seconds.
\begin{lstlisting}[numbers=none]
.venv/bin/python scripts/download_data.py --dataset tokyo
\end{lstlisting}

\noindent\textbf{20261105T052634.748331Z}: exit 0, 61.876 seconds.
\begin{lstlisting}[numbers=none]
.venv/bin/python scripts/run_originals.py --limit 90
\end{lstlisting}

\noindent\textbf{20261105T052649.085364Z}: exit 0, 108.436 seconds.
\begin{lstlisting}[numbers=none]
.venv/bin/python scripts/download_models.py
\end{lstlisting}

\noindent\textbf{20261105T052655.188712Z}: exit 0, 31.692 seconds.
\begin{lstlisting}[numbers=none]
.venv/bin/python -c 'import cv2,json,platform,torch,psutil,importlib.metadata as m;from pathlib import Path;v=cv2.VideoCapture('"'"'video.mp4'"'"');d=dict(platform=platform.platform(),machine=platform.machine(),python=platform.python_version(),cpu=platform.processor(),physical_memory_bytes=psutil.virtual_memory().total,torch_mps_available=torch.backends.mps.is_available(),torch_cuda_available=torch.cuda.is_available(),packages={p:m.version(p) for p in ['"'"'numpy'"'"','"'"'opencv-contrib-python'"'"','"'"'torch'"'"','"'"'torchvision'"'"','"'"'kornia'"'"','"'"'scipy'"'"','"'"'matplotlib'"'"','"'"'Pillow'"'"','"'"'onnx'"'"','"'"'onnxruntime'"'"']},video=dict(width=v.get(3),height=v.get(4),fps=v.get(5),frames=v.get(7)));v.release();Path('"'"'results/environment.json'"'"').write_text(json.dumps(d,indent=2));print(json.dumps(d,indent=2))'
\end{lstlisting}

\noindent\textbf{20261105T052736.685439Z}: exit 0, 4.905 seconds.
\begin{lstlisting}[numbers=none]
.venv/bin/python -m pytest -q
\end{lstlisting}

\noindent\textbf{20261105T052950.790466Z}: exit 0, 2.418 seconds.
\begin{lstlisting}[numbers=none]
.venv/bin/python scripts/download_data.py --dataset hpatches
\end{lstlisting}

\noindent\textbf{20261105T053026.926415Z}: exit 0, 5.494 seconds.
\begin{lstlisting}[numbers=none]
.venv/bin/python -c 'from stitch.features import LocalEngine;from stitch.io import read_image;from stitch.retrieval import GlobalEngine;import torch,time;torch.set_num_threads(2);im=read_image('"'"'reference.jpg'"'"')[0];e=LocalEngine('"'"'superpoint'"'"',512);a=e.extract(im);start=time.perf_counter();p,q,ids=e.match(a,a);print('"'"'SP points'"'"',len(a.points),'"'"'matches'"'"',len(p),'"'"'matching_ms'"'"',(time.perf_counter()-start)*1000);g=GlobalEngine();print('"'"'CosPlace dimension'"'"',g.extract(im).shape)'
\end{lstlisting}

\noindent\textbf{20261105T053131.878032Z}: exit 1, 0.067 seconds.
\begin{lstlisting}[numbers=none]
.venv/bin/python -m compileall -q src scripts tests
\end{lstlisting}

\noindent\textbf{20261105T053145.439997Z}: exit 0, 0.170 seconds.
\begin{lstlisting}[numbers=none]
.venv/bin/python -m compileall -q src scripts tests
\end{lstlisting}

\noindent\textbf{20261105T053145.639983Z}: exit 0, 2.712 seconds.
\begin{lstlisting}[numbers=none]
.venv/bin/python scripts/download_data.py --dataset hpatches
\end{lstlisting}

\noindent\textbf{20261105T053148.401823Z}: exit 0, 2.216 seconds.
\begin{lstlisting}[numbers=none]
.venv/bin/python scripts/prepare_data.py
\end{lstlisting}

\noindent\textbf{20261105T053203.373611Z}: exit 0, 968.671 seconds.
\begin{lstlisting}[numbers=none]
.venv/bin/python scripts/evaluate.py --learned-exhaustive
\end{lstlisting}

\noindent\textbf{20261105T053231.430339Z}: exit 0, 25.342 seconds.
\begin{lstlisting}[numbers=none]
.venv/bin/python scripts/train.py --steps 120
\end{lstlisting}

\noindent\textbf{20261105T053256.841737Z}: exit 0, 26.273 seconds.
\begin{lstlisting}[numbers=none]
.venv/bin/python scripts/export_models.py
\end{lstlisting}

\noindent\textbf{20261105T053428.648382Z}: exit 0, 2.723 seconds.
\begin{lstlisting}[numbers=none]
.venv/bin/python -m pytest -q
\end{lstlisting}

\noindent\textbf{20261105T053431.416021Z}: exit 0, 0.891 seconds.
\begin{lstlisting}[numbers=none]
.venv/bin/python -m pip freeze --all
\end{lstlisting}

\noindent\textbf{20261105T053452.374058Z}: exit 0, 1.796 seconds.
\begin{lstlisting}[numbers=none]
.venv/bin/python -c 'import cv2,json,numpy as np;from pathlib import Path;rows=json.loads(Path('"'"'data/tokyo/selected.json'"'"').read_text());panels=[];first=rows[0]['"'"'place_id'"'"'];selected=[r for r in rows if r['"'"'place_id'"'"']==first];print(selected);[(print(r['"'"'image'"'"'],r['"'"'condition'"'"'],float(cv2.imread('"'"'data/tokyo/'"'"'+r['"'"'image'"'"']).mean()))) for r in selected];panels=[cv2.resize(cv2.imread('"'"'data/tokyo/'"'"'+r['"'"'image'"'"']),(320,240)) for r in selected];cv2.imwrite('"'"'figures/tokyo_inspection.jpg'"'"',np.hstack(panels))'
\end{lstlisting}

\noindent\textbf{20261105T053456.339375Z}: exit 0, 0.562 seconds.
\begin{lstlisting}[numbers=none]
curl -fsSL --connect-timeout 10 --max-time 30 -o results/sources/aachen_readme.md https://data.ciirc.cvut.cz/public/projects/2020VisualLocalization/Aachen-Day-Night/README_Aachen-Day-Night_v1_1.md
\end{lstlisting}

\noindent\textbf{20261105T053456.950016Z}: exit 0, 0.749 seconds.
\begin{lstlisting}[numbers=none]
curl -fsSL --connect-timeout 10 --max-time 30 -o results/sources/robotcar_readme.md https://data.ciirc.cvut.cz/public/projects/2020VisualLocalization/RobotCar-Seasons/README_RobotCar-Seasons.md
\end{lstlisting}

\noindent\textbf{20261105T053457.777340Z}: exit 0, 0.739 seconds.
\begin{lstlisting}[numbers=none]
curl -fsSL --connect-timeout 10 --max-time 30 -o results/sources/cmu_readme.md https://data.ciirc.cvut.cz/public/projects/2020VisualLocalization/Extended-CMU-Seasons/README_Extended-CMU-Seasons.md
\end{lstlisting}

\noindent\textbf{20261105T053458.549454Z}: exit 0, 0.401 seconds.
\begin{lstlisting}[numbers=none]
curl -fsSL --connect-timeout 10 --max-time 30 -o results/sources/hpatches_readme.md https://raw.githubusercontent.com/hpatches/hpatches-dataset/master/README.md
\end{lstlisting}

\noindent\textbf{20261105T054650.922828Z}: exit 0, 9.309 seconds.
\begin{lstlisting}[numbers=none]
.venv/bin/python tep.py
\end{lstlisting}

\noindent\textbf{20261105T054652.047516Z}: exit 0, 0.038 seconds.
\begin{lstlisting}[numbers=none]
.venv/bin/python -c 'import json;from pathlib import Path;d=json.loads(Path('"'"'results/training.json'"'"').read_text());print([(r['"'"'protocol'"'"'],r['"'"'variant'"'"'],r['"'"'augmentation'"'"'],r['"'"'seed'"'"'],round(r['"'"'recall1'"'"'],3)) for r in d['"'"'results'"'"'] if r['"'"'protocol'"'"'] in ['"'"'tokyo_day_to_day'"'"','"'"'hpatches_viewpoint'"'"']])'
\end{lstlisting}

\noindent\textbf{20261105T054700.276524Z}: exit 0, 5.598 seconds.
\begin{lstlisting}[numbers=none]
.venv/bin/stitch index data/places.json
\end{lstlisting}

\noindent\textbf{20261105T054705.910108Z}: exit 0, 28.593 seconds.
\begin{lstlisting}[numbers=none]
.venv/bin/stitch query data/places.json test.jpg results/raw/cli_query.png
\end{lstlisting}

\noindent\textbf{20261105T054742.415068Z}: exit 0, 16.932 seconds.
\begin{lstlisting}[numbers=none]
.venv/bin/python scripts/probe_data.py
\end{lstlisting}

\noindent\textbf{20261105T054759.395798Z}: exit 0, 2.074 seconds.
\begin{lstlisting}[numbers=none]
.venv/bin/python -c 'import sys,json;from pathlib import Path;sys.path.insert(0,'"'"'scripts'"'"');from download_data import member,TOKYO;rows=json.loads(Path('"'"'data/tokyo_zip_directory.json'"'"').read_text());i=next(r for r in rows if r['"'"'name'"'"'].endswith('"'"'/Readme.txt'"'"'));member(i,TOKYO,Path.cwd()/'"'"'results/sources/tokyo_readme.txt'"'"')'
\end{lstlisting}

\noindent\textbf{20261105T054953.631155Z}: exit 1, 2.152 seconds.
\begin{lstlisting}[numbers=none]
.venv/bin/python scripts/qualitative.py
\end{lstlisting}

\noindent\textbf{20261105T054954.748109Z}: exit 0, 0.081 seconds.
\begin{lstlisting}[numbers=none]
.venv/bin/python -c 'import csv,json;from pathlib import Path;rs=list(csv.DictReader(Path('"'"'results/retrieval.csv'"'"').open()));print('"'"'\n'"'"'.join(str(r) for r in rs if r['"'"'protocol'"'"'].startswith('"'"'tokyo_'"'"') and r['"'"'method'"'"'] in ['"'"'A_SIFT'"'"','"'"'B_SP_LightGlue'"'"','"'"'C_CosPlace'"'"','"'"'D_full_top3'"'"']));p=json.loads(Path('"'"'results/pairs.json'"'"').read_text());import numpy as np;print([(s,m,np.mean([r['"'"'homography_success'"'"'] for r in p if r['"'"'protocol'"'"']==s and r['"'"'method'"'"']==m]),np.mean([r['"'"'accepted'"'"'] for r in p if r['"'"'protocol'"'"']==s and r['"'"'method'"'"']==m])) for s in ['"'"'hpatches_illumination'"'"','"'"'hpatches_viewpoint'"'"'] for m in ['"'"'sift'"'"','"'"'superpoint'"'"']])'
\end{lstlisting}

\noindent\textbf{20261105T054955.815465Z}: exit 1, 1.022 seconds.
\begin{lstlisting}[numbers=none]
.venv/bin/python scripts/demo.py
\end{lstlisting}

\noindent\textbf{20261105T054956.868306Z}: exit 1, 1.010 seconds.
\begin{lstlisting}[numbers=none]
.venv/bin/python scripts/profile.py
\end{lstlisting}

\noindent\textbf{20261105T054957.912795Z}: exit 0, 118.743 seconds.
\begin{lstlisting}[numbers=none]
.venv/bin/python scripts/generate_artifacts.py
\end{lstlisting}

\noindent\textbf{20261105T055242.908276Z}: exit 0, 11.955 seconds.
\begin{lstlisting}[numbers=none]
.venv/bin/python scripts/qualitative.py
\end{lstlisting}

\noindent\textbf{20261105T055254.934633Z}: exit 0, 103.781 seconds.
\begin{lstlisting}[numbers=none]
.venv/bin/python scripts/demo.py
\end{lstlisting}

\noindent\textbf{20261105T055438.792760Z}: exit 0, 152.640 seconds.
\begin{lstlisting}[numbers=none]
.venv/bin/python scripts/benchmark_runtime.py
\end{lstlisting}

\noindent\textbf{20261105T055441.193860Z}: exit 0, 10.173 seconds.
\begin{lstlisting}[numbers=none]
/Users/bhattacharjee/Desktop/internship/Stitch-research/.venv/bin/python scripts/benchmark_runtime.py --method classical
\end{lstlisting}

\noindent\textbf{20261105T055451.421071Z}: exit 0, 84.772 seconds.
\begin{lstlisting}[numbers=none]
/Users/bhattacharjee/Desktop/internship/Stitch-research/.venv/bin/python scripts/benchmark_runtime.py --method local
\end{lstlisting}

\noindent\textbf{20261105T055455.482303Z}: exit 0, 0.045 seconds.
\begin{lstlisting}[numbers=none]
.venv/bin/python -c 'import json,csv;from pathlib import Path;p=json.loads(Path('"'"'results/predictions.json'"'"').read_text());print('"'"'Open-set'"'"',[(m,sum(r['"'"'accepted'"'"'] for r in p if r['"'"'protocol'"'"']=='"'"'open_set'"'"' and r['"'"'method'"'"']==m)) for m in ['"'"'A_SIFT'"'"','"'"'B_SP_LightGlue'"'"','"'"'D_full_top1'"'"','"'"'D_full_top3'"'"','"'"'D_full_top5'"'"','"'"'D_no_geometry'"'"']]);rows=list(csv.DictReader(Path('"'"'results/retrieval.csv'"'"').open()));print('"'"'Synthetic severe'"'"',[r for r in rows if r['"'"'protocol'"'"']=='"'"'synthetic_night_severe'"'"' and r['"'"'method'"'"'] in ['"'"'A_SIFT'"'"','"'"'B_SP_LightGlue'"'"','"'"'C_CosPlace'"'"','"'"'D_full_top3'"'"']]);print('"'"'Ablation'"'"',[r for r in rows if r['"'"'protocol'"'"']=='"'"'tokyo_day_to_day'"'"'])'
\end{lstlisting}

\noindent\textbf{20261105T055616.354179Z}: exit 0, 6.809 seconds.
\begin{lstlisting}[numbers=none]
/Users/bhattacharjee/Desktop/internship/Stitch-research/.venv/bin/python scripts/benchmark_runtime.py --method retrieval
\end{lstlisting}

\noindent\textbf{20261105T055623.199031Z}: exit 0, 48.002 seconds.
\begin{lstlisting}[numbers=none]
/Users/bhattacharjee/Desktop/internship/Stitch-research/.venv/bin/python scripts/benchmark_runtime.py --method full
\end{lstlisting}

\noindent\textbf{20261105T055711.499630Z}: exit 0, 15.141 seconds.
\begin{lstlisting}[numbers=none]
.venv/bin/python scripts/generate_artifacts.py
\end{lstlisting}

\noindent\textbf{20261105T060521.591059Z}: exit 0, 0.166 seconds.
\begin{lstlisting}[numbers=none]
.venv/bin/python -c 'from pathlib import Path;import subprocess;Path('"'"'docs'"'"').mkdir(exist_ok=True);Path('"'"'docs/UPSTREAM_README.md'"'"').write_bytes(subprocess.check_output(['"'"'git'"'"','"'"'show'"'"','"'"'HEAD:README.md'"'"']))'
\end{lstlisting}

\noindent\textbf{20261105T060521.796123Z}: exit 0, 2.502 seconds.
\begin{lstlisting}[numbers=none]
.venv/bin/python -m pytest -q
\end{lstlisting}

\noindent\textbf{20261105T060524.343033Z}: exit 0, 39.402 seconds.
\begin{lstlisting}[numbers=none]
.venv/bin/stitch video data/places.json video.mp4 results/raw/cli_video.mp4 --limit 8
\end{lstlisting}

\noindent\textbf{20261105T060527.912335Z}: exit 0, 16.320 seconds.
\begin{lstlisting}[numbers=none]
bash scripts/setup.sh
\end{lstlisting}

\noindent\textbf{20261105T060528.169115Z}: exit 0, 1.718 seconds.
\begin{lstlisting}[numbers=none]
.venv/bin/python -m pip install --upgrade pip
\end{lstlisting}

\noindent\textbf{20261105T060529.934447Z}: exit 0, 0.542 seconds.
\begin{lstlisting}[numbers=none]
.venv/bin/python -m pip install -r requirements.txt
\end{lstlisting}

\noindent\textbf{20261105T060530.633366Z}: exit 0, 3.969 seconds.
\begin{lstlisting}[numbers=none]
.venv/bin/python -m pip install -e . --no-deps
\end{lstlisting}

\noindent\textbf{20261105T060534.649113Z}: exit 0, 3.841 seconds.
\begin{lstlisting}[numbers=none]
.venv/bin/python -m pip install -e vendor/LightGlue --no-deps
\end{lstlisting}

\noindent\textbf{20261105T060538.608904Z}: exit 0, 5.351 seconds.
\begin{lstlisting}[numbers=none]
.venv/bin/python scripts/download_models.py
\end{lstlisting}

\noindent\textbf{20261105T060544.023207Z}: exit 0, 0.205 seconds.
\begin{lstlisting}[numbers=none]
.venv/bin/python scripts/snapshot_environment.py
\end{lstlisting}

\noindent\textbf{20261105T060603.812353Z}: exit 0, 4.141 seconds.
\begin{lstlisting}[numbers=none]
make data
\end{lstlisting}

\noindent\textbf{20261105T060604.722962Z}: exit 0, 1.164 seconds.
\begin{lstlisting}[numbers=none]
.venv/bin/python scripts/download_data.py
\end{lstlisting}

\noindent\textbf{20261105T060605.920926Z}: exit 0, 2.028 seconds.
\begin{lstlisting}[numbers=none]
.venv/bin/python scripts/prepare_data.py
\end{lstlisting}

\noindent\textbf{20261105T060752.805044Z}: exit 2, 23.068 seconds.
\begin{lstlisting}[numbers=none]
make thesis
\end{lstlisting}

\noindent\textbf{20261105T060752.894846Z}: in progress at artifact generation.
\begin{lstlisting}[numbers=none]
/Library/Developer/CommandLineTools/usr/bin/make -C thesis thesis
\end{lstlisting}

\noindent\textbf{20261105T060858.429350Z}: exit 2, 19.297 seconds.
\begin{lstlisting}[numbers=none]
make thesis
\end{lstlisting}

\noindent\textbf{20261105T060858.510106Z}: exit 2, 19.212 seconds.
\begin{lstlisting}[numbers=none]
/Library/Developer/CommandLineTools/usr/bin/make -C thesis thesis
\end{lstlisting}

\noindent\textbf{20261105T060927.786710Z}: exit 0, 1.320 seconds.
\begin{lstlisting}[numbers=none]
.venv/bin/python scripts/download_data.py --verify-only
\end{lstlisting}

\noindent\textbf{20261105T060929.156015Z}: exit 0, 0.040 seconds.
\begin{lstlisting}[numbers=none]
.venv/bin/python -c 'import shutil;print({p:shutil.which(p) for p in ['"'"'gs'"'"','"'"'pdftoppm'"'"','"'"'mutool'"'"','"'"'latexmk'"'"','"'"'bibtex'"'"','"'"'makeindex'"'"']})'
\end{lstlisting}

\noindent\textbf{20261105T061002.417428Z}: exit 0, 0.212 seconds.
\begin{lstlisting}[numbers=none]
make clean
\end{lstlisting}

\noindent\textbf{20261105T061002.667976Z}: exit 2, 30.203 seconds.
\begin{lstlisting}[numbers=none]
make thesis
\end{lstlisting}

\noindent\textbf{20261105T061002.720547Z}: exit 2, 30.144 seconds.
\begin{lstlisting}[numbers=none]
/Library/Developer/CommandLineTools/usr/bin/make -C thesis thesis
\end{lstlisting}

\noindent\textbf{20261105T061119.642113Z}: exit 0, 0.958 seconds.
\begin{lstlisting}[numbers=none]
.venv/bin/python scripts/verify_thesis.py
\end{lstlisting}

\noindent\textbf{20261105T061120.635443Z}: exit 0, 5.394 seconds.
\begin{lstlisting}[numbers=none]
.venv/bin/python -c 'from pypdf import PdfReader;r=PdfReader('"'"'thesis/thesis.pdf'"'"');print([(i+1,p.extract_text()[:95].replace('"'"'\n'"'"','"'"' '"'"')) for i,p in enumerate(r.pages) if '"'"'Table 8.1:'"'"' in p.extract_text() or '"'"'Chapter 1'"'"' in p.extract_text()[:100] or '"'"'Bibliography'"'"' in p.extract_text()[:100]]);print(r.outline[:3])'
\end{lstlisting}

\noindent\textbf{20261105T061153.880326Z}: exit 0, 0.892 seconds.
\begin{lstlisting}[numbers=none]
pdftoppm -f 67 -l 67 -scale-to 1300 -png -singlefile thesis/thesis.pdf results/raw/thesis_table_check
\end{lstlisting}

\noindent\textbf{20261105T061154.809827Z}: exit 0, 0.202 seconds.
\begin{lstlisting}[numbers=none]
pdftoppm -f 37 -l 37 -scale-to 1300 -png -singlefile thesis/thesis.pdf results/raw/thesis_architecture_check
\end{lstlisting}

\noindent\textbf{20261105T061155.045666Z}: exit 0, 0.237 seconds.
\begin{lstlisting}[numbers=none]
pdftoppm -f 78 -l 78 -scale-to 1300 -png -singlefile thesis/thesis.pdf results/raw/thesis_overlay_check
\end{lstlisting}

\noindent\textbf{20261105T061155.320668Z}: exit 0, 0.093 seconds.
\begin{lstlisting}[numbers=none]
pdftoppm -f 1 -l 1 -scale-to 1300 -png -singlefile thesis/thesis.pdf results/raw/thesis_title_check
\end{lstlisting}

\noindent\textbf{20261105T061322.200930Z}: exit 0, 0.145 seconds.
\begin{lstlisting}[numbers=none]
make clean
\end{lstlisting}

\noindent\textbf{20261105T061322.384509Z}: exit 0, 28.754 seconds.
\begin{lstlisting}[numbers=none]
make thesis
\end{lstlisting}

\noindent\textbf{20261105T061322.430691Z}: exit 0, 28.703 seconds.
\begin{lstlisting}[numbers=none]
/Library/Developer/CommandLineTools/usr/bin/make -C thesis thesis
\end{lstlisting}

\noindent\textbf{20261105T061616.989356Z}: exit 0, 109.893 seconds.
\begin{lstlisting}[numbers=none]
.venv/bin/python scripts/demo.py
\end{lstlisting}

\noindent\textbf{20261105T061630.344476Z}: exit 0, 0.151 seconds.
\begin{lstlisting}[numbers=none]
.venv/bin/python -c 'from pathlib import Path; import json; p=Path('"'"'results/demo.json'"'"'); d=json.loads(p.read_text()); print('"'"'Demo input:'"'"', *d['"'"'source_images'"'"']); print('"'"'Output:'"'"', Path('"'"'demo/comparison.mp4'"'"').resolve()); print('"'"'Video write in progress:'"'"', Path('"'"'demo/comparison.mp4'"'"').stat().st_size, '"'"'bytes'"'"')'
\end{lstlisting}

\noindent\textbf{20261105T061818.031121Z}: exit 0, 0.379 seconds.
\begin{lstlisting}[numbers=none]
.venv/bin/python -c 'import cv2; c=cv2.VideoCapture('"'"'demo/comparison.mp4'"'"'); n=0; ok,f=c.read(); assert ok, '"'"'Video cannot be decoded'"'"'; shape=f.shape; n=1; ok,f=c.read();
while ok:
 n+=1; ok,f=c.read()
c.release(); assert n==48, n; print('"'"'Verified fresh demo:'"'"', n, '"'"'frames;'"'"', shape[1], '"'"'x'"'"', shape[0])'
\end{lstlisting}

\noindent\textbf{20261105T061818.441559Z}: exit 1, 0.126 seconds.
\begin{lstlisting}[numbers=none]
/usr/bin/open demo/comparison.mp4
\end{lstlisting}

\noindent\textbf{20261105T061832.531575Z}: exit 0, 0.261 seconds.
\begin{lstlisting}[numbers=none]
/usr/bin/open -a 'QuickTime Player' demo/comparison.mp4
\end{lstlisting}

\noindent\textbf{20261105T062500.396651Z}: exit 0, 0.438 seconds.
\begin{lstlisting}[numbers=none]
.venv/bin/python -c 'from pypdf import PdfReader;from pathlib import Path;files=[Path('"'"'../PRONOB-BHATTACHARJEE-latex/PRONOB BHATTACHARJEE.pdf'"'"'),Path('"'"'../PRONOB-BHATTACHARJEE-Overleaf/main.pdf'"'"'),Path('"'"'../IITR_template.pdf'"'"')];
for path in files:
 r=PdfReader(path);print('"'"'\nFILE'"'"',path,'"'"'PAGES'"'"',len(r.pages),'"'"'SIZE'"'"',r.pages[0].mediabox);print('"'"'\n'"'"'.join((p.extract_text() or '"'"''"'"')[:6500] for p in r.pages[:6]))'
\end{lstlisting}

\noindent\textbf{20261105T062547.765248Z}: exit 0, 0.623 seconds.
\begin{lstlisting}[numbers=none]
.venv/bin/python -c 'from pathlib import Path;import zipfile;from PIL import Image;from pypdf import PdfReader;print('"'"'ASSETS'"'"');
for p in Path('"'"'..'"'"').glob('"'"'*'"'"'):
 if p.suffix.lower() in ['"'"'.png'"'"','"'"'.jpg'"'"']:print(p.name,Image.open(p).size)
for p in Path('"'"'..'"'"').glob('"'"'*.zip'"'"'):
 z=zipfile.ZipFile(p);print('"'"'ARCHIVE'"'"',p.name,len(z.namelist()));print([n for n in z.namelist() if n.endswith(('"'"'.sty'"'"','"'"'.cls'"'"','"'"'.map'"'"'))][:30])
r=PdfReader('"'"'../PRONOB-BHATTACHARJEE-Overleaf/main.pdf'"'"');print('"'"'THESIS PAGE TITLES'"'"');print([(i+1,(r.pages[i].extract_text() or '"'"''"'"')[:180].replace('"'"'\n'"'"','"'"' '"'"')) for i in [7,8,9,10,11,12,13,14,15,20,35]]);print('"'"'PAGE1 IMAGES'"'"',[(v.name,len(v.data),v.image.size) for v in r.pages[0].images]);print('"'"'TEXT POSITIONS'"'"');
r.pages[0].extract_text(visitor_text=lambda t,cm,tm,f,s:print(repr(t[:85]),'"'"'xy'"'"',tm[4:6],'"'"'size'"'"',s) if t.strip() else None)'
\end{lstlisting}

\noindent\textbf{20261105T062548.450208Z}: exit 0, 0.200 seconds.
\begin{lstlisting}[numbers=none]
pdftoppm -f 1 -l 1 -scale-to 1300 -png -singlefile ../PRONOB-BHATTACHARJEE-Overleaf/main.pdf results/raw/reference_report_cover
\end{lstlisting}

\noindent\textbf{20261105T062548.680032Z}: exit 0, 0.219 seconds.
\begin{lstlisting}[numbers=none]
pdftoppm -f 14 -l 14 -scale-to 1300 -png -singlefile ../PRONOB-BHATTACHARJEE-Overleaf/main.pdf results/raw/reference_report_body
\end{lstlisting}

\noindent\textbf{20261105T062548.928585Z}: exit 0, 0.874 seconds.
\begin{lstlisting}[numbers=none]
pdftoppm -f 1 -l 1 -scale-to 1500 -png -singlefile ../IITR_template.pdf results/raw/reference_slide_title
\end{lstlisting}

\noindent\textbf{20261105T062549.832160Z}: exit 0, 1.149 seconds.
\begin{lstlisting}[numbers=none]
pdftoppm -f 4 -l 4 -scale-to 1500 -png -singlefile ../IITR_template.pdf results/raw/reference_slide_body
\end{lstlisting}

\noindent\textbf{20261105T062654.091555Z}: exit 0, 0.232 seconds.
\begin{lstlisting}[numbers=none]
.venv/bin/python -c 'from pypdf import PdfReader; r=PdfReader('"'"'../PRONOB-BHATTACHARJEE-Overleaf/main.pdf'"'"');
for i in [15,16]:
 print('"'"'PAGE'"'"',i+1);r.pages[i].extract_text(visitor_text=lambda t,cm,tm,f,s:print(repr(t[:70]),tm[4:6],s) if t.strip() and (tm[5]>590 or tm[5]<90) else None)
print('"'"'PDF modules'"'"');import importlib.util;print({n:bool(importlib.util.find_spec(n)) for n in ['"'"'pptx'"'"','"'"'fitz'"'"','"'"'reportlab'"'"']})'
\end{lstlisting}

\noindent\textbf{20261105T062706.126122Z}: exit 0, 0.214 seconds.
\begin{lstlisting}[numbers=none]
.venv/bin/python -c 'from pypdf import PdfReader;r=PdfReader('"'"'../PRONOB-BHATTACHARJEE-Overleaf/thesis-graphics.pdf'"'"');
for i in [0,1]:
 print('"'"'GRAPHICS'"'"',i+1);r.pages[i].extract_text(visitor_operand_before=lambda op,args,cm,tm:print(op,args) if op in [b'"'"'cm'"'"',b'"'"'Do'"'"'] else None)
r=PdfReader('"'"'../PRONOB-BHATTACHARJEE-Overleaf/main.pdf'"'"');r.pages[15].extract_text(visitor_text=lambda t,cm,tm,f,s:print(repr(t[:60]),tm[4:6],s) if t.strip() else None)'
\end{lstlisting}

\noindent\textbf{20261105T063014.121353Z}: exit 0, 2.735 seconds.
\begin{lstlisting}[numbers=none]
.venv/bin/python -m pip install python-pptx==1.0.2
\end{lstlisting}

\noindent\textbf{20261105T063015.406488Z}: exit 0, 0.276 seconds.
\begin{lstlisting}[numbers=none]
.venv/bin/python scripts/prepare_iitr_assets.py
\end{lstlisting}

\noindent\textbf{20261105T063234.202515Z}: exit 0, 0.033 seconds.
\begin{lstlisting}[numbers=none]
.venv/bin/python scripts/restyle_iitr_report.py
\end{lstlisting}

\noindent\textbf{20261105T063234.266409Z}: exit 0, 25.767 seconds.
\begin{lstlisting}[numbers=none]
make thesis
\end{lstlisting}

\noindent\textbf{20261105T063234.342166Z}: exit 0, 25.687 seconds.
\begin{lstlisting}[numbers=none]
/Library/Developer/CommandLineTools/usr/bin/make -C thesis thesis
\end{lstlisting}

\noindent\textbf{20261105T063640.164469Z}: exit 0, 0.560 seconds.
\begin{lstlisting}[numbers=none]
.venv/bin/python scripts/generate_presentation_data.py
\end{lstlisting}

\noindent\textbf{20261105T063640.756639Z}: exit 2, 6.796 seconds.
\begin{lstlisting}[numbers=none]
make -C presentation slides
\end{lstlisting}

\noindent\textbf{20261105T063748.190935Z}: exit 0, 0.200 seconds.
\begin{lstlisting}[numbers=none]
pdftoppm -f 1 -l 1 -scale-to 1300 -png -singlefile thesis/thesis.pdf results/raw/iitr_report_cover
\end{lstlisting}

\noindent\textbf{20261105T063748.424041Z}: exit 0, 0.130 seconds.
\begin{lstlisting}[numbers=none]
pdftoppm -f 2 -l 2 -scale-to 1300 -png -singlefile thesis/thesis.pdf results/raw/iitr_report_inner_cover
\end{lstlisting}

\noindent\textbf{20261105T063748.584871Z}: exit 0, 0.143 seconds.
\begin{lstlisting}[numbers=none]
pdftoppm -f 4 -l 4 -scale-to 1300 -png -singlefile thesis/thesis.pdf results/raw/iitr_report_declaration
\end{lstlisting}

\noindent\textbf{20261105T063748.758544Z}: exit 0, 0.191 seconds.
\begin{lstlisting}[numbers=none]
pdftoppm -f 20 -l 20 -scale-to 1300 -png -singlefile thesis/thesis.pdf results/raw/iitr_report_chapter
\end{lstlisting}

\noindent\textbf{20261105T063826.276948Z}: exit 0, 25.947 seconds.
\begin{lstlisting}[numbers=none]
make -C presentation slides
\end{lstlisting}

\noindent\textbf{20261105T063852.288769Z}: exit 0, 1.574 seconds.
\begin{lstlisting}[numbers=none]
pdftoppm -f 1 -l 2 -scale-to 1400 -png presentation/Stitch_IITR_Presentation.pdf results/raw/iitr_slide
\end{lstlisting}

\noindent\textbf{20261105T063853.897155Z}: exit 0, 0.357 seconds.
\begin{lstlisting}[numbers=none]
pdftoppm -f 15 -l 15 -scale-to 1400 -png -singlefile presentation/Stitch_IITR_Presentation.pdf results/raw/iitr_slide_results
\end{lstlisting}

\noindent\textbf{20261105T063906.533229Z}: exit 0, 0.215 seconds.
\begin{lstlisting}[numbers=none]
.venv/bin/python -c 'from pypdf import PdfReader;r=PdfReader('"'"'thesis/thesis.pdf'"'"');
for i in [0,1]:
 print('"'"'PAGE'"'"',i+1,r.pages[i].mediabox)
 def show(t,cm,tm,font,size):
  if t.strip():print(repr(t[:85]),'"'"'position'"'"',(round(tm[4]*cm[0]+tm[5]*cm[2]+cm[4],2),round(tm[4]*cm[1]+tm[5]*cm[3]+cm[5],2)))
 r.pages[i].extract_text(visitor_text=show)'
\end{lstlisting}

\noindent\textbf{20261105T063906.841090Z}: exit 0, 0.371 seconds.
\begin{lstlisting}[numbers=none]
pdftoppm -f 16 -l 16 -scale-to 1300 -png -singlefile ../PRONOB-BHATTACHARJEE-Overleaf/main.pdf results/raw/reference_report_chapter
\end{lstlisting}

\noindent\textbf{20261105T064000.220624Z}: exit 0, 0.308 seconds.
\begin{lstlisting}[numbers=none]
.venv/bin/python -c 'from pypdf import PdfReader;
for path,i in [('"'"'../PRONOB-BHATTACHARJEE-Overleaf/main.pdf'"'"',15),('"'"'thesis/thesis.pdf'"'"',19)]:
 print(path);r=PdfReader(path)
 def show(t,cm,tm,f,s):
  if t.strip() and tm[5]>490:print(repr(t[:65]),round(tm[4],3),round(tm[5],3),s)
 r.pages[i].extract_text(visitor_text=show)'
\end{lstlisting}

\noindent\textbf{20261105T064016.391671Z}: exit 0, 0.330 seconds.
\begin{lstlisting}[numbers=none]
.venv/bin/python -c 'from pypdf import PdfReader;
for path,i in [('"'"'../PRONOB-BHATTACHARJEE-Overleaf/main.pdf'"'"',16),('"'"'thesis/thesis.pdf'"'"',20)]:
 print(path);rows=[];PdfReader(path).pages[i].extract_text(visitor_text=lambda t,cm,tm,f,s:rows.append((t[:65],tm[4],tm[5],s)) if t.strip() else None);print(rows[:5]);print(rows[-2:])'
\end{lstlisting}

\noindent\textbf{20261105T064112.752468Z}: exit 0, 0.217 seconds.
\begin{lstlisting}[numbers=none]
.venv/bin/python -c 'from pypdf import PdfReader;r=PdfReader('"'"'presentation/Stitch_IITR_Presentation.pdf'"'"');print(r.pages[1].extract_text());
def f(t,cm,tm,font,size):
 if '"'"'INDIAN'"'"' in t:print(repr(t),cm,tm,size)
r.pages[1].extract_text(visitor_text=f)'
\end{lstlisting}

\noindent\textbf{20261105T064113.000883Z}: exit 0, 24.164 seconds.
\begin{lstlisting}[numbers=none]
make thesis
\end{lstlisting}

\noindent\textbf{20261105T064113.045480Z}: exit 0, 24.114 seconds.
\begin{lstlisting}[numbers=none]
/Library/Developer/CommandLineTools/usr/bin/make -C thesis thesis
\end{lstlisting}

\noindent\textbf{20261105T064404.078519Z}: exit 0, 0.461 seconds.
\begin{lstlisting}[numbers=none]
.venv/bin/python scripts/verify_iitr_documents.py
\end{lstlisting}

\noindent\textbf{20261105T064404.596654Z}: exit 0, 0.468 seconds.
\begin{lstlisting}[numbers=none]
.venv/bin/python scripts/package_iitr_documents.py
\end{lstlisting}

\noindent\textbf{20261105T064447.228720Z}: exit 0, 16.910 seconds.
\begin{lstlisting}[numbers=none]
.venv/bin/python scripts/check_document_packages.py
\end{lstlisting}

\noindent\textbf{20261105T064448.337993Z}: exit 0, 0.711 seconds.
\begin{lstlisting}[numbers=none]
pdftoppm -f 2 -l 2 -scale-to 1400 -png -singlefile presentation/Stitch_IITR_Presentation.pdf results/raw/iitr_slide_author_final
\end{lstlisting}

\noindent\textbf{20261105T064449.080706Z}: exit 0, 0.141 seconds.
\begin{lstlisting}[numbers=none]
.venv/bin/python -c 'from pypdf import PdfReader;r=PdfReader('"'"'thesis/thesis.pdf'"'"');p=r.pages[19];out=[];p.extract_text(visitor_text=lambda t,cm,tm,f,s:out.append((t[:50],tm[4],tm[5],s)) if t.strip() and tm[5]>490 else None);print(out);print('"'"'body bounds:'"'"',r.pages[19].mediabox)'
\end{lstlisting}

\noindent\textbf{20261105T064504.180674Z}: exit 0, 0.178 seconds.
\begin{lstlisting}[numbers=none]
.venv/bin/python scripts/snapshot_environment.py
\end{lstlisting}

\noindent\textbf{20261105T064603.528739Z}: exit 0, 1.219 seconds.
\begin{lstlisting}[numbers=none]
.venv/bin/python scripts/verify_iitr_documents.py
\end{lstlisting}

\noindent\textbf{20261105T064604.779912Z}: exit 0, 1.132 seconds.
\begin{lstlisting}[numbers=none]
.venv/bin/python scripts/verify_thesis.py
\end{lstlisting}

\noindent\textbf{20261105T064605.947423Z}: exit 0, 0.096 seconds.
\begin{lstlisting}[numbers=none]
.venv/bin/python -c 'from pathlib import Path;from pypdf import PdfReader;import json;d=Path('"'"'deliverables'"'"');print('"'"'\n'"'"'.join(str(p.resolve())+'"'"' | '"'"'+str(p.stat().st_size)+'"'"' bytes'"'"' for p in sorted(d.iterdir()) if p.is_file()));print(Path('"'"'results/standalone_packages_verification.json'"'"').read_text())'
\end{lstlisting}

\noindent\textbf{20261105T064814.639158Z}: exit 0, 2.171 seconds.
\begin{lstlisting}[numbers=none]
.venv/bin/python -m pytest -q
\end{lstlisting}

\noindent\textbf{20261105T064922.537044Z}: exit 1, 7.880 seconds.
\begin{lstlisting}[numbers=none]
.venv/bin/python -m pip download --no-deps --dest .cache/wheels opencv-python==4.11.0.86
\end{lstlisting}

\noindent\textbf{20261105T064934.773406Z}: exit 0, 6.122 seconds.
\begin{lstlisting}[numbers=none]
.venv/bin/python -m pip download --no-deps --dest .cache/wheels opencv-python==4.11.0.86
\end{lstlisting}

\noindent\textbf{20261105T065059.529486Z}: exit 0, 0.330 seconds.
\begin{lstlisting}[numbers=none]
.venv/bin/python -m pip uninstall -y opencv-contrib-python
\end{lstlisting}

\noindent\textbf{20261105T065059.892978Z}: exit 0, 0.437 seconds.
\begin{lstlisting}[numbers=none]
.venv/bin/python -m pip install --no-deps .cache/wheels/opencv_python-4.11.0.86-cp37-abi3-macosx_13_0_arm64.whl
\end{lstlisting}

\noindent\textbf{20261105T065100.361237Z}: exit 0, 0.268 seconds.
\begin{lstlisting}[numbers=none]
.venv/bin/python -m pip check
\end{lstlisting}

\noindent\textbf{20261105T065147.072171Z}: exit 0, 22.727 seconds.
\begin{lstlisting}[numbers=none]
.venv/bin/python -m pytest -q
\end{lstlisting}

\noindent\textbf{20261105T065209.841463Z}: exit 0, 0.062 seconds.
\begin{lstlisting}[numbers=none]
.venv/bin/python -m compileall -q src scripts tests
\end{lstlisting}

\noindent\textbf{20261105T065209.932280Z}: exit 0, 0.159 seconds.
\begin{lstlisting}[numbers=none]
.venv/bin/python scripts/snapshot_environment.py
\end{lstlisting}

\noindent\textbf{20261105T065258.722422Z}: exit 1, 1.701 seconds.
\begin{lstlisting}[numbers=none]
.venv/bin/python scripts/verify_research.py
\end{lstlisting}

\noindent\textbf{20261105T065300.468906Z}: exit 0, 3.647 seconds.
\begin{lstlisting}[numbers=none]
.venv/bin/stitch index data/places.json
\end{lstlisting}

\noindent\textbf{20261105T065304.162382Z}: exit 0, 20.794 seconds.
\begin{lstlisting}[numbers=none]
.venv/bin/stitch query data/places.json test.jpg results/raw/audit_cli_query.png
\end{lstlisting}

\noindent\textbf{20261105T065325.002021Z}: exit 0, 20.568 seconds.
\begin{lstlisting}[numbers=none]
.venv/bin/stitch video data/places.json video.mp4 results/raw/audit_cli_video.mp4 --limit 2
\end{lstlisting}

\noindent\textbf{20261105T065353.751749Z}: exit 0, 6.728 seconds.
\begin{lstlisting}[numbers=none]
.venv/bin/python scripts/verify_research.py
\end{lstlisting}

\noindent\textbf{20261105T065433.490353Z}: exit 0, 42.477 seconds.
\begin{lstlisting}[numbers=none]
make package
\end{lstlisting}

\noindent\textbf{20261105T065433.536594Z}: exit 0, 21.449 seconds.
\begin{lstlisting}[numbers=none]
/Library/Developer/CommandLineTools/usr/bin/make -C thesis thesis
\end{lstlisting}

\noindent\textbf{20261105T065434.593113Z}: exit 0, 0.248 seconds.
\begin{lstlisting}[numbers=none]
.venv/bin/python -
\end{lstlisting}

\noindent\textbf{20261105T065455.024877Z}: exit 0, 0.220 seconds.
\begin{lstlisting}[numbers=none]
.venv/bin/python scripts/prepare_iitr_assets.py
\end{lstlisting}

\noindent\textbf{20261105T065455.276702Z}: exit 0, 0.134 seconds.
\begin{lstlisting}[numbers=none]
.venv/bin/python scripts/generate_presentation_data.py
\end{lstlisting}

\noindent\textbf{20261105T065455.441088Z}: exit 0, 19.009 seconds.
\begin{lstlisting}[numbers=none]
/Library/Developer/CommandLineTools/usr/bin/make -C presentation slides
\end{lstlisting}

\noindent\textbf{20261105T065514.494629Z}: exit 0, 1.006 seconds.
\begin{lstlisting}[numbers=none]
.venv/bin/python scripts/verify_iitr_documents.py
\end{lstlisting}

\noindent\textbf{20261105T065515.531302Z}: exit 0, 0.432 seconds.
\begin{lstlisting}[numbers=none]
.venv/bin/python scripts/package_iitr_documents.py
\end{lstlisting}

\noindent\textbf{20261105T065543.483631Z}: exit 0, 1.004 seconds.
\begin{lstlisting}[numbers=none]
.venv/bin/python -
\end{lstlisting}

\noindent\textbf{20261105T065550.691204Z}: exit 0, 27.870 seconds.
\begin{lstlisting}[numbers=none]
make check-documents
\end{lstlisting}

\noindent\textbf{20261105T065550.745068Z}: exit 0, 8.323 seconds.
\begin{lstlisting}[numbers=none]
.venv/bin/python scripts/verify_research.py
\end{lstlisting}

\noindent\textbf{20261105T065559.119409Z}: exit 0, 0.883 seconds.
\begin{lstlisting}[numbers=none]
.venv/bin/python scripts/verify_thesis.py
\end{lstlisting}

\noindent\textbf{20261105T065600.035437Z}: exit 0, 1.063 seconds.
\begin{lstlisting}[numbers=none]
.venv/bin/python scripts/verify_iitr_documents.py
\end{lstlisting}

\noindent\textbf{20261105T065601.130152Z}: exit 0, 17.426 seconds.
\begin{lstlisting}[numbers=none]
.venv/bin/python scripts/check_document_packages.py
\end{lstlisting}

\noindent\textbf{20261105T065713.174447Z}: exit 1, 0.040 seconds.
\begin{lstlisting}[numbers=none]
.venv/bin/python -
\end{lstlisting}

\noindent\textbf{20261105T065724.753702Z}: exit 0, 17.102 seconds.
\begin{lstlisting}[numbers=none]
.venv/bin/python scripts/check_document_packages.py
\end{lstlisting}

\noindent\textbf{20261105T065748.296746Z}: exit 0, 0.056 seconds.
\begin{lstlisting}[numbers=none]
.venv/bin/python -m compileall -q src scripts tests
\end{lstlisting}

\noindent\textbf{20261105T065748.435570Z}: exit 0, 0.324 seconds.
\begin{lstlisting}[numbers=none]
.venv/bin/python -m pip check
\end{lstlisting}

\noindent\textbf{20261105T171849.404819Z}: exit 0, 0.872 seconds.
\begin{lstlisting}[numbers=none]
.venv/bin/python -
\end{lstlisting}

\noindent\textbf{20261105T171917.747993Z}: exit 0, 0.231 seconds.
\begin{lstlisting}[numbers=none]
.venv/bin/python -
\end{lstlisting}

\noindent\textbf{20261105T172020.452180Z}: exit 0, 3.390 seconds.
\begin{lstlisting}[numbers=none]
.venv/bin/python -m pytest -q
\end{lstlisting}

\noindent\textbf{20261105T172023.874238Z}: exit 0, 6.582 seconds.
\begin{lstlisting}[numbers=none]
.venv/bin/python scripts/verify_research.py
\end{lstlisting}

\noindent\textbf{20261105T172109.500695Z}: exit 0, 42.914 seconds.
\begin{lstlisting}[numbers=none]
make -j2 PYTHON=.venv/bin/python package
\end{lstlisting}

\noindent\textbf{20261105T172109.551511Z}: exit 0, 22.013 seconds.
\begin{lstlisting}[numbers=none]
/Library/Developer/CommandLineTools/usr/bin/make -C thesis PYTHON=/Users/bhattacharjee/Desktop/internship/Stitch-research/.venv/bin/python thesis
\end{lstlisting}

\noindent\textbf{20261105T172131.599762Z}: exit 0, 0.179 seconds.
\begin{lstlisting}[numbers=none]
.venv/bin/python scripts/prepare_iitr_assets.py
\end{lstlisting}

\noindent\textbf{20261105T172131.807537Z}: exit 0, 0.124 seconds.
\begin{lstlisting}[numbers=none]
.venv/bin/python scripts/generate_presentation_data.py
\end{lstlisting}

\noindent\textbf{20261105T172131.964820Z}: exit 0, 18.968 seconds.
\begin{lstlisting}[numbers=none]
/Library/Developer/CommandLineTools/usr/bin/make -C presentation PYTHON=/Users/bhattacharjee/Desktop/internship/Stitch-research/.venv/bin/python slides
\end{lstlisting}

\noindent\textbf{20261105T172150.974376Z}: exit 0, 0.999 seconds.
\begin{lstlisting}[numbers=none]
.venv/bin/python scripts/verify_iitr_documents.py
\end{lstlisting}

\noindent\textbf{20261105T172152.003575Z}: exit 0, 0.408 seconds.
\begin{lstlisting}[numbers=none]
.venv/bin/python scripts/package_iitr_documents.py
\end{lstlisting}

\noindent\textbf{20261105T172222.235375Z}: exit 0, 17.423 seconds.
\begin{lstlisting}[numbers=none]
.venv/bin/python scripts/verify_model_artifacts.py
\end{lstlisting}

\noindent\textbf{20261105T172250.066037Z}: exit 0, 47.782 seconds.
\begin{lstlisting}[numbers=none]
make check-documents
\end{lstlisting}

\noindent\textbf{20261105T172250.181407Z}: exit 0, 11.736 seconds.
\begin{lstlisting}[numbers=none]
.venv/bin/python scripts/verify_research.py
\end{lstlisting}

\noindent\textbf{20261105T172301.957670Z}: exit 0, 17.216 seconds.
\begin{lstlisting}[numbers=none]
.venv/bin/python scripts/verify_model_artifacts.py
\end{lstlisting}

\noindent\textbf{20261105T172319.221315Z}: exit 0, 0.755 seconds.
\begin{lstlisting}[numbers=none]
.venv/bin/python scripts/verify_thesis.py
\end{lstlisting}

\noindent\textbf{20261105T172320.006398Z}: exit 0, 1.075 seconds.
\begin{lstlisting}[numbers=none]
.venv/bin/python scripts/verify_iitr_documents.py
\end{lstlisting}

\noindent\textbf{20261105T172321.118316Z}: exit 0, 16.726 seconds.
\begin{lstlisting}[numbers=none]
.venv/bin/python scripts/check_document_packages.py
\end{lstlisting}

\noindent\textbf{20261105T172324.634750Z}: exit 0, 0.548 seconds.
\begin{lstlisting}[numbers=none]
.venv/bin/python -
\end{lstlisting}

\noindent\textbf{20261105T172325.214992Z}: exit 0, 0.050 seconds.
\begin{lstlisting}[numbers=none]
.venv/bin/python -m compileall -q src scripts tests
\end{lstlisting}

\noindent\textbf{20261105T172325.294683Z}: exit 0, 0.308 seconds.
\begin{lstlisting}[numbers=none]
.venv/bin/python -m pip check
\end{lstlisting}

\noindent\textbf{20261105T172354.394797Z}: exit 0, 21.212 seconds.
\begin{lstlisting}[numbers=none]
.venv/bin/stitch query data/places.json test.jpg results/raw/second_audit_cli_query.png
\end{lstlisting}

\noindent\textbf{20261105T172415.653597Z}: exit 0, 0.262 seconds.
\begin{lstlisting}[numbers=none]
.venv/bin/python scripts/snapshot_environment.py
\end{lstlisting}

\noindent\textbf{20261105T205001.296062Z}: exit 0, 51.011 seconds.
\begin{lstlisting}[numbers=none]
make thesis
\end{lstlisting}

\noindent\textbf{20261105T205001.384604Z}: exit 0, 50.918 seconds.
\begin{lstlisting}[numbers=none]
/Library/Developer/CommandLineTools/usr/bin/make -C thesis PYTHON=/Users/bhattacharjee/Desktop/internship/Stitch-research/.venv/bin/python thesis
\end{lstlisting}

\noindent\textbf{20261105T205157.650794Z}: exit 0, 3.916 seconds.
\begin{lstlisting}[numbers=none]
.venv/bin/python -
\end{lstlisting}

\noindent\textbf{20261105T205201.608392Z}: exit 0, 3.081 seconds.
\begin{lstlisting}[numbers=none]
.venv/bin/python -m pytest -q
\end{lstlisting}

\noindent\textbf{20261105T205204.722463Z}: exit 0, 67.738 seconds.
\begin{lstlisting}[numbers=none]
make package
\end{lstlisting}

\noindent\textbf{20261105T205204.817739Z}: exit 0, 45.750 seconds.
\begin{lstlisting}[numbers=none]
/Library/Developer/CommandLineTools/usr/bin/make -C thesis PYTHON=/Users/bhattacharjee/Desktop/internship/Stitch-research/.venv/bin/python thesis
\end{lstlisting}

\noindent\textbf{20261105T205250.602065Z}: exit 0, 0.191 seconds.
\begin{lstlisting}[numbers=none]
.venv/bin/python scripts/prepare_iitr_assets.py
\end{lstlisting}

\noindent\textbf{20261105T205250.822360Z}: exit 0, 0.128 seconds.
\begin{lstlisting}[numbers=none]
.venv/bin/python scripts/generate_presentation_data.py
\end{lstlisting}

\noindent\textbf{20261105T205250.982915Z}: exit 0, 19.021 seconds.
\begin{lstlisting}[numbers=none]
/Library/Developer/CommandLineTools/usr/bin/make -C presentation PYTHON=/Users/bhattacharjee/Desktop/internship/Stitch-research/.venv/bin/python slides
\end{lstlisting}

\noindent\textbf{20261105T205310.047707Z}: exit 0, 1.185 seconds.
\begin{lstlisting}[numbers=none]
.venv/bin/python scripts/verify_iitr_documents.py
\end{lstlisting}

\noindent\textbf{20261105T205311.262414Z}: exit 0, 1.195 seconds.
\begin{lstlisting}[numbers=none]
.venv/bin/python scripts/package_iitr_documents.py
\end{lstlisting}

\noindent\textbf{20261105T205404.307368Z}: exit 0, 0.616 seconds.
\begin{lstlisting}[numbers=none]
.venv/bin/python -c 'from scripts.generate_extended_evidence import write_manifest; result=write_manifest(); print("Evidence assets hashed:",len(result["generated_outputs"]))'
\end{lstlisting}

\noindent\textbf{20261105T205404.954443Z}: exit 0, 56.487 seconds.
\begin{lstlisting}[numbers=none]
make check-documents
\end{lstlisting}

\noindent\textbf{20261105T205405.046440Z}: exit 0, 12.064 seconds.
\begin{lstlisting}[numbers=none]
.venv/bin/python scripts/verify_research.py
\end{lstlisting}

\noindent\textbf{20261105T205417.155986Z}: exit 0, 17.374 seconds.
\begin{lstlisting}[numbers=none]
.venv/bin/python scripts/verify_model_artifacts.py
\end{lstlisting}

\noindent\textbf{20261105T205434.576635Z}: exit 0, 4.240 seconds.
\begin{lstlisting}[numbers=none]
.venv/bin/python scripts/verify_thesis.py
\end{lstlisting}

\noindent\textbf{20261105T205438.847011Z}: exit 0, 1.198 seconds.
\begin{lstlisting}[numbers=none]
.venv/bin/python scripts/verify_iitr_documents.py
\end{lstlisting}

\noindent\textbf{20261105T205440.074142Z}: exit 0, 21.363 seconds.
\begin{lstlisting}[numbers=none]
.venv/bin/python scripts/check_document_packages.py
\end{lstlisting}

\noindent\textbf{20261105T205542.620927Z}: exit 0, 5.155 seconds.
\begin{lstlisting}[numbers=none]
latexmk -pdf -interaction=nonstopmode -halt-on-error -outdir=build main.tex
\end{lstlisting}

\noindent\textbf{20261105T205543.738644Z}: exit 0, 0.062 seconds.
\begin{lstlisting}[numbers=none]
.venv/bin/python -m compileall -q src scripts tests
\end{lstlisting}

\noindent\textbf{20261105T205543.831580Z}: exit 0, 0.338 seconds.
\begin{lstlisting}[numbers=none]
.venv/bin/python -m pip check
\end{lstlisting}

\noindent\textbf{20261105T205552.631065Z}: exit 0, 4.266 seconds.
\begin{lstlisting}[numbers=none]
.venv/bin/python scripts/verify_thesis.py
\end{lstlisting}

\noindent\textbf{20261105T205556.927358Z}: exit 0, 1.196 seconds.
\begin{lstlisting}[numbers=none]
.venv/bin/python scripts/verify_iitr_documents.py
\end{lstlisting}

\noindent\textbf{20261105T205558.153498Z}: exit 0, 1.159 seconds.
\begin{lstlisting}[numbers=none]
.venv/bin/python scripts/package_iitr_documents.py
\end{lstlisting}

\noindent\textbf{20261105T205559.347193Z}: exit 0, 21.175 seconds.
\begin{lstlisting}[numbers=none]
.venv/bin/python scripts/check_document_packages.py
\end{lstlisting}

\noindent\textbf{20261105T205650.800177Z}: exit 0, 4.185 seconds.
\begin{lstlisting}[numbers=none]
.venv/bin/python -
\end{lstlisting}

\noindent\textbf{20261107T063039.615283Z}: exit 0, 0.646 seconds.
\begin{lstlisting}[numbers=none]
.venv/bin/python -m pip install --upgrade pip
\end{lstlisting}

\noindent\textbf{20261107T063040.292584Z}: exit 0, 0.353 seconds.
\begin{lstlisting}[numbers=none]
.venv/bin/python -m pip install -r requirements.txt
\end{lstlisting}

\noindent\textbf{20261107T063040.704200Z}: exit 0, 1.951 seconds.
\begin{lstlisting}[numbers=none]
.venv/bin/python -m pip install -e . --no-deps
\end{lstlisting}

\noindent\textbf{20261107T063042.686558Z}: exit 0, 1.805 seconds.
\begin{lstlisting}[numbers=none]
.venv/bin/python -m pip install -e vendor/LightGlue --no-deps
\end{lstlisting}

\noindent\textbf{20261107T063044.524714Z}: exit 0, 3.830 seconds.
\begin{lstlisting}[numbers=none]
.venv/bin/python scripts/download_models.py
\end{lstlisting}

\noindent\textbf{20261107T063048.387226Z}: exit 0, 0.141 seconds.
\begin{lstlisting}[numbers=none]
.venv/bin/python scripts/snapshot_environment.py
\end{lstlisting}

\noindent\textbf{20261107T063048.568397Z}: exit 0, 1.207 seconds.
\begin{lstlisting}[numbers=none]
.venv/bin/python scripts/download_data.py
\end{lstlisting}

\noindent\textbf{20261107T063049.813980Z}: exit 0, 2.008 seconds.
\begin{lstlisting}[numbers=none]
.venv/bin/python scripts/prepare_data.py
\end{lstlisting}

\noindent\textbf{20261107T103959.273112Z}: exit 1, 2.785 seconds.
\begin{lstlisting}[numbers=none]
.venv/bin/python -m pytest -q
\end{lstlisting}

\noindent\textbf{20261107T104017.187591Z}: exit 0, 2.768 seconds.
\begin{lstlisting}[numbers=none]
.venv/bin/python -m pytest -q
\end{lstlisting}

\noindent\textbf{20261107T104030.546778Z}: exit 0, 12.277 seconds.
\begin{lstlisting}[numbers=none]
.venv/bin/python scripts/run_originals.py --limit 90
\end{lstlisting}

\noindent\textbf{20261107T104042.888982Z}: in progress at artifact generation.
\begin{lstlisting}[numbers=none]
.venv/bin/python scripts/evaluate.py --learned-exhaustive
\end{lstlisting}

\noindent\textbf{20261107T183120.484098Z}: exit 0, 0.695 seconds.
\begin{lstlisting}[numbers=none]
.venv/bin/python -m pip install --upgrade pip
\end{lstlisting}

\noindent\textbf{20261107T183121.210974Z}: exit 0, 0.341 seconds.
\begin{lstlisting}[numbers=none]
.venv/bin/python -m pip install -r requirements.txt
\end{lstlisting}

\noindent\textbf{20261107T183121.610703Z}: exit 0, 2.047 seconds.
\begin{lstlisting}[numbers=none]
.venv/bin/python -m pip install -e . --no-deps
\end{lstlisting}

\noindent\textbf{20261107T183123.688979Z}: exit 0, 1.844 seconds.
\begin{lstlisting}[numbers=none]
.venv/bin/python -m pip install -e vendor/LightGlue --no-deps
\end{lstlisting}

\noindent\textbf{20261107T183125.566290Z}: exit 0, 3.745 seconds.
\begin{lstlisting}[numbers=none]
.venv/bin/python scripts/download_models.py
\end{lstlisting}

\noindent\textbf{20261107T183129.341507Z}: exit 0, 0.136 seconds.
\begin{lstlisting}[numbers=none]
.venv/bin/python scripts/snapshot_environment.py
\end{lstlisting}

\noindent\textbf{20261107T195701.490617Z}: exit 0, 7.960 seconds.
\begin{lstlisting}[numbers=none]
.venv/bin/python -m pip install --upgrade pip
\end{lstlisting}

\noindent\textbf{20261107T195709.485124Z}: exit 0, 0.380 seconds.
\begin{lstlisting}[numbers=none]
.venv/bin/python -m pip install -r requirements.txt
\end{lstlisting}

\noindent\textbf{20261107T195709.917630Z}: exit 1, 7.951 seconds.
\begin{lstlisting}[numbers=none]
.venv/bin/python -m pip install -e . --no-deps
\end{lstlisting}

\noindent\textbf{20261107T195736.805577Z}: exit 0, 0.526 seconds.
\begin{lstlisting}[numbers=none]
.venv/bin/python -m pip install --upgrade pip
\end{lstlisting}

\noindent\textbf{20261107T195737.363098Z}: exit 0, 0.326 seconds.
\begin{lstlisting}[numbers=none]
.venv/bin/python -m pip install -r requirements.txt
\end{lstlisting}

\noindent\textbf{20261107T195737.739104Z}: exit 0, 1.933 seconds.
\begin{lstlisting}[numbers=none]
.venv/bin/python -m pip install -e . --no-deps
\end{lstlisting}

\noindent\textbf{20261107T195739.703592Z}: exit 0, 1.778 seconds.
\begin{lstlisting}[numbers=none]
.venv/bin/python -m pip install -e vendor/LightGlue --no-deps
\end{lstlisting}

\noindent\textbf{20261107T195741.512731Z}: exit 0, 2.967 seconds.
\begin{lstlisting}[numbers=none]
.venv/bin/python scripts/download_models.py
\end{lstlisting}

\noindent\textbf{20261107T195744.510862Z}: exit 0, 0.137 seconds.
\begin{lstlisting}[numbers=none]
.venv/bin/python scripts/snapshot_environment.py
\end{lstlisting}

\noindent\textbf{20261109T073045.677679Z}: exit 0, 0.004 seconds.
\begin{lstlisting}[numbers=none]
/Library/Developer/CommandLineTools/usr/bin/make -C report PYTHON=/Users/bhattacharjee/Desktop/internship/Stitch-research/.venv/bin/python thesis
\end{lstlisting}

\noindent\textbf{20261109T074313.788133Z}: exit 0, 3.131 seconds.
\begin{lstlisting}[numbers=none]
.venv/bin/python -m pytest -q
\end{lstlisting}

\noindent\textbf{20261109T074323.382473Z}: exit 0, 1.300 seconds.
\begin{lstlisting}[numbers=none]
.venv/bin/python scripts/download_data.py
\end{lstlisting}

\noindent\textbf{20261109T074324.712374Z}: exit 0, 1.613 seconds.
\begin{lstlisting}[numbers=none]
.venv/bin/python scripts/prepare_data.py
\end{lstlisting}

\noindent\textbf{20261109T074326.362324Z}: exit 0, 11.225 seconds.
\begin{lstlisting}[numbers=none]
.venv/bin/python scripts/run_originals.py --limit 90
\end{lstlisting}

\noindent\textbf{20261109T074350.495586Z}: exit 0, 893.074 seconds.
\begin{lstlisting}[numbers=none]
.venv/bin/python scripts/evaluate.py --learned-exhaustive
\end{lstlisting}

\noindent\textbf{20261109T075918.479055Z}: exit 0, 22.127 seconds.
\begin{lstlisting}[numbers=none]
.venv/bin/python scripts/train.py --steps 120
\end{lstlisting}

\noindent\textbf{20261109T075940.750069Z}: exit 0, 21.902 seconds.
\begin{lstlisting}[numbers=none]
.venv/bin/python scripts/export_models.py
\end{lstlisting}

\noindent\textbf{20261109T080010.888464Z}: exit 0, 101.332 seconds.
\begin{lstlisting}[numbers=none]
.venv/bin/python scripts/demo.py
\end{lstlisting}

\noindent\textbf{20261109T080152.338377Z}: exit 0, 164.440 seconds.
\begin{lstlisting}[numbers=none]
.venv/bin/python scripts/benchmark_runtime.py
\end{lstlisting}

\noindent\textbf{20261109T080154.286035Z}: exit 0, 9.125 seconds.
\begin{lstlisting}[numbers=none]
/Users/bhattacharjee/Desktop/internship/Stitch-research/.venv/bin/python scripts/benchmark_runtime.py --method classical
\end{lstlisting}

\noindent\textbf{20261109T080203.449039Z}: exit 0, 87.552 seconds.
\begin{lstlisting}[numbers=none]
/Users/bhattacharjee/Desktop/internship/Stitch-research/.venv/bin/python scripts/benchmark_runtime.py --method local
\end{lstlisting}

\noindent\textbf{20261109T080331.147238Z}: exit 0, 7.356 seconds.
\begin{lstlisting}[numbers=none]
/Users/bhattacharjee/Desktop/internship/Stitch-research/.venv/bin/python scripts/benchmark_runtime.py --method retrieval
\end{lstlisting}

\noindent\textbf{20261109T080338.541730Z}: exit 0, 57.976 seconds.
\begin{lstlisting}[numbers=none]
/Users/bhattacharjee/Desktop/internship/Stitch-research/.venv/bin/python scripts/benchmark_runtime.py --method full
\end{lstlisting}

\noindent\textbf{20261109T080447.349271Z}: in progress at artifact generation.
\begin{lstlisting}[numbers=none]
/Library/Developer/CommandLineTools/usr/bin/make -C report PYTHON=/Users/bhattacharjee/Desktop/internship/Stitch-research/.venv/bin/python thesis
\end{lstlisting}

\fi
